

\section{Cellular structures on cyclotomic $q$-W-Schur algebras}
\label{sec:cellular}

In this section, we establish the cellular structure for $\wS_\bfu(n)$ via the corresponding cellular structure of the cyclotomic $q$-Schur category. We also study the cell filtration of the restriction and induction of the cell modules of $\wS_\bfu$.

    
 \subsection{Cellular bases and cell modules}   

\begin{definition} [\cite{SW24Schur, SSW25}]
\label{def:qSchur}
    The affine $q$-Schur category  $\qASch$ is a strict $\kk$-linear monoidal category with generating objects $a\in \Z_{\ge 1}$ and $u\in \kk$. We denote the object $a\in \N$ by $\stra$ and denote the object $u \in \kk$ by a red strand labeled by $u$ as $\stru .$
The morphisms are generated by 
\begin{equation} 
\label{generator-affschur}
  \begin{tikzpicture}[anchorbase, scale=0.4, color=\clr] 
                \draw[-,line width=1pt] (-2,0.2) to (-1.5,1);
	\draw[-,line width=1pt ] (-1,0.2) to (-1.5,1);
	\draw[-,line width=1.5pt] (-1.5,1) to (-1.5,1.8);
                \draw (-1,0) node{$\scriptstyle {b}$};
                \draw (-2,0) node{$\scriptstyle {a}$};
                \draw (-1.5,2.2) node{$\scriptstyle {a+b}$};
                \end{tikzpicture}, 
                \quad 
  \begin{tikzpicture}[anchorbase, scale=0.4, color=\clr]
                \draw[-,line width=1pt] (-2,1.8) to (-1.5,1);
	\draw[-,line width=1pt ] (-1,1.8) to (-1.5,1);
	\draw[-,line width=1.5pt] (-1.5,1) to (-1.5,0.2);
                \draw (-1,2.2) node{$\scriptstyle {b}$};
                \draw (-2,2.2) node{$\scriptstyle {a}$};
                \draw (-1.5,0) node{$\scriptstyle {a+b}$};
                \end{tikzpicture}, 
    \quad  \;
    \crossingpos , 
        \quad \;
    ~\xdota ,
 \end{equation}
and
\begin{equation}
\label{redcrossing}
\rightcrossing 
\;\; (\text{traverse-up}), 
                \qquad
                \leftcrossing 
 \;\; (\text{traverse-down}),    
\end{equation}
subject to relations \eqref{webassoc}--\eqref{intergralballon} for the generators in \eqref{generator-affschur} and additional relations \eqref{redslider}--\eqref{redbraid} involving also the traverse-ups and traverse-downs \eqref{redcrossing}:
%
\begin{align}
\label{redslider}
\begin{tikzpicture}[anchorbase,scale=.6,color=\clr]
	\draw[-,line width=1pt,color=\cred] (0.4,0) to (-0.6,1);
  \draw (-.6,1.2) node{$\scriptstyle \red{u}$};
	\draw[-,line width=1pt] (0.08,0) to (0.08,1);
	\draw[-,line width=1pt] (0.1,0) to (0.1,.6) to (.5,1);
        \node at (0.6,1.18) {$\scriptstyle a$};
        \node at (0.1,1.2) {$\scriptstyle b$};
\end{tikzpicture}
& =
\begin{tikzpicture}[anchorbase,scale=.6,color=\clr]
	\draw[-,line width=1pt,color=\cred] (0.7,0) to (-0.3,1);
  \draw (-.3,1.16) node{$\scriptstyle \red{u}$};
	\draw[-,line width=1pt] (0.08,0) to (0.08,1);
	\draw[-,line width=1pt] (0.1,0) to (0.1,.2) to (.9,1);
        \node at (0.9,1.18) {$\scriptstyle a$};
        \node at (0.1,1.2) {$\scriptstyle b$};
\end{tikzpicture},
\quad
\begin{tikzpicture}[anchorbase,scale=.6,color=\clr]
	\draw[-,line width=1pt,color=\cred] (-0.4,0) to (0.6,1);
  \draw (.6,1.16) node{$\scriptstyle \red{u}$};
	\draw[-,line width=1pt] (-0.08,0) to (-0.08,1);
	\draw[-,line width=1pt] (-0.1,0) to (-0.1,.6) to (-.5,1);
         \node at (-0.1,1.18) {$\scriptstyle b$};
        \node at (-0.6,1.2) {$\scriptstyle a$};
\end{tikzpicture}
\!\!=\!\!
\begin{tikzpicture}[anchorbase,scale=.6,color=\clr]
	\draw[-,line width=1pt,color=\cred] (-0.7,0) to (0.3,1);
 \draw (.3,1.16) node{$\scriptstyle \red{u}$};
	\draw[-,line width=1pt] (-0.08,0) to (-0.08,1);
	\draw[-,line width=1pt] (-0.1,0) to (-0.1,.2) to (-.9,1);
         \node at (-0.1,1.2) {$\scriptstyle b$};
        \node at (-0.95,1.18) {$\scriptstyle a$};
\end{tikzpicture}, 
\quad
\begin{tikzpicture}[baseline=-3.3mm,scale=.6,color=\clr]
\draw[-,line width=1pt,color=\cred] (0.4,.2) to (-0.6,-.8);
\draw (-.6,-1.05) node{$\scriptstyle \red{u}$};
\draw[-,line width=1pt] (0.08,0.2) to (0.08,-.75);
\draw[-,line width=1pt] (0.1,0.2) to (0.1,-.4) to (.5,-.8);
\node at (0.6,-1.05) {$\scriptstyle c$};
\node at (0.07,-1.05) {$\scriptstyle b$};
\end{tikzpicture}
\!\!=\!\!
\begin{tikzpicture}[baseline=-3.3mm,scale=.6,color=\clr]
\draw[-,line width=1pt,color=\cred] (0.7,0.2) to (-0.3,-.8);
\draw (-.3,-1.0) node{$\scriptstyle \red{u}$};
\draw[-,line width=1pt] (0.08,0.2) to (0.08,-.75);
\draw[-,line width=1pt] (0.1,0.2) to (0.1,0) to (.9,-.8);
\node at (1,-1.0) {$\scriptstyle c$};
\node at (0.1,-1.0) {$\scriptstyle b$};
\end{tikzpicture} ,
\quad
\begin{tikzpicture}[baseline=-3.3mm,scale=.6,color=\clr]
\draw[-,line width=1pt,color=\cred] (-0.4,0.2) to (0.6,-.8);
\draw (.6,-.91) node{$\scriptstyle \red{u}$};
\draw[-,line width=1pt] (-0.08,0.2) to (-0.08,-.75);
\draw[-,line width=1pt] (-0.1,0.2) to (-0.1,-.4) to (-.5,-.8);
\node at (-0.1,-1.0) {$\scriptstyle b$};
\node at (-0.6,-1.0) {$\scriptstyle a$};
\end{tikzpicture}
\!\!=\!\!
\begin{tikzpicture}[baseline=-3.3mm,scale=.6,color=\clr]
\draw[-,line width=1pt,color=\cred] (-0.7,0.2) to (0.3,-.8);
\draw (.3,-1) node{$\scriptstyle \red{u}$};
\draw[-,line width=1pt] (-0.08,0.2) to (-0.08,-.75);
\draw[-,line width=1pt] (-0.1,0.2) to (-0.1,0) to (-.9,-.8);
\node at (-0.1,-1.0) {$\scriptstyle b$};
\node at (-0.95,-1.0) {$\scriptstyle a$};
\end{tikzpicture} ,
\\
\label{redcross2}
\begin{split}
\begin{tikzpicture}[baseline = 7pt, scale=0.3, color=\clr]
\draw[-,line width=1pt,color=\cred](0.1,-.22)to (0,2.5);
\draw (0.1,-.5) node{$\scriptstyle \red{u}$};
\draw[-,line width=1.2pt](-.5,-.1) to[out=45, in=down] (.5,1.4);
\draw[-,line width=1.2pt](.5,1.4) to[out=up, in=300] (-.5,2.4);
\draw (-.5,-0.5) node{$\scriptstyle {a}$};
\end{tikzpicture}
 & =
\sum_{0\le t \le a} (-1)^t q^{-t(a-t)} u^t
\begin{tikzpicture}[baseline = 7pt, scale=0.3, color=\clr]
\draw[-,line width=1pt,color=\cred](-.5,-.25)to (-.5,2.15);
\draw (-.5,-.6) node{$\scriptstyle \red{u}$};
\draw[-,line width=1pt] (-1.3,2.15) to (-1.3,-0.2);   
\draw (-1.3,1) \bdot;
\draw (-2.5,1.4) node{$\scriptstyle \omega_{a-t}$};
\draw (-1.3,-.6) node{$\scriptstyle {a}$};
\end{tikzpicture} ~, 
\\    
\begin{tikzpicture}[baseline = 7pt, scale=0.3, color=\clr]
\draw[-,line width=1pt,color=\cred](-0.1,-.6)to (-0.1,2.2);
\draw (-0.1,-.9) node{$\scriptstyle \red{u}$};
\draw[-,line width=1.2pt](.5,-.5) to[out=135, in=down] (-.5,1);
\draw[-,line width=1.2pt](-.5,1) to[out=up, in=270] (.5,2.2);
\draw (.5,-0.9) node{$\scriptstyle {a}$};
\end{tikzpicture}
&=
\sum_{0\le t \le a} (-1)^t q^{-t(a-t)} u^t
\begin{tikzpicture}[baseline = 7pt, scale=0.3, color=\clr]
\draw[-,line width=1pt,color=\cred](-1.8,-.25)to (-1.8,2.15);
\draw (-1.8,-.6) node{$\scriptstyle \red{u}$};
\draw[-,line width=1pt] (-1,2.15) to (-1,-0.25);          
\draw (-1,1) \bdot;
\draw (0.4,1)  node{$\scriptstyle \omega_{a-t}$};
\draw (-1,-.6) node{$\scriptstyle {a}$};
\end{tikzpicture} ,
\end{split}
\\
 \label{redbraid}
\begin{tikzpicture}[anchorbase, scale=0.35, color=\clr]
\draw[-,line width=1pt](0,2) to (1.5,-.35);
\draw[-,line width=4pt, white](0,-.35) to (1.5,2); to (1.5,-.35);
\draw[-,line width=1pt] (0,-.35) to (1.5,2);
\draw[-,line width=1pt,color=\cred](.8,2) to[out=down,in=90] (0.2,1) to [out=down,in=up] (.8,-.4);
\draw (.8,-.8) node{$\scriptstyle \red{u}$};
\node at (0, -.7) {$\scriptstyle a$};
\node at (1.5, -.7) {$\scriptstyle b$};
\end{tikzpicture}
& = 
\sum_{s=0}^{\min (a,b)} (-1)^s q^{s(s-1)/2} (q^{-1} -q)^s [s]_q! 
\begin{tikzpicture}[anchorbase, scale=0.35, color=\clr]
\draw[-,line width=1.5pt](0,-.35) to   (0,0);
\draw[-,line width=1.5pt](0,1.65) to   (0,2);
\draw[-,line width=1.5pt](1.5,-.35) to   (1.5,0);
\draw[-,line width=1.5pt](1.5,1.65) to   (1.5,2);
\draw[-,thick](0,0) to   (0,1.65);
\draw[-,thick](1.5,0) to   (1.5,1.65);
\draw[-,line width=1pt](0,1.65) to   (1.5,0);
\draw[-,line width=4pt, white] (0.2,0.2) to (1.3,1.35);
\draw[-,line width=1pt] (0,0) to (1.5,1.65);
\draw[-,line width=1pt,color=\cred](.8,2) to
[out=-60,in=60]
(.8,-.4);
\draw (.8,-.7) node{$\scriptstyle \red{u}$};
\node at (0, -.7) {$\scriptstyle a$};
\node at (1.5, -.7) {$\scriptstyle b$};
\node at (-0.3, 1.1) {$\scriptstyle s$};
\node at (1.8, 1.1) {$\scriptstyle s$};
\draw (1.5, .7) \bdot;
\end{tikzpicture} . 
\end{align}
\end{definition}

Let $\Lambda_{\text{st}}^{1+\ell}(n)$ be the set of all strict $(1+\ell)$-multicompositions of $n$. Write 
\[
\Lambda_{\text{st}}^{1+\ell}:=\bigcup_{m\in \N}\Lambda_{\text{st}}^{1+\ell}(m).
\]
Thus, the set of objects of $\qASch$
can be identified with $ \bigcup_{\ell\in \N} \big(\Lambda_{\text{st}}^{1+\ell}\times \kk^{\ell} \big)$.

Below we will identify 
$\lambda^{(0)}\red{u_1}\ldots \red{u_\ell} \lambda^{(\ell)}$ with $\lambda=(\lambda^{(0)},\ldots, \lambda^{(\ell)})$ once $\bfu$ is fixed in \eqref{equ:parau}.
For any fixed $\bfu$, the cyclotomic $q$-Schur category is the quotient of the $\kk$-linear category $\qSchu$
    by the relations 
    \begin{equation}
    \label{eq:unitlambda}
       \unit{\lambda}=0 \text{ for all } \lambda =(\la^{(0)},\la^{(1)},\ldots,\la^{(\ell)}) \in \Lambda_{\text{st}}^{1+\ell} \text{ with }
\lambda^{(0)}\neq \emptyset.
\end{equation}

By \eqref{eq:unitlambda}, the set of objects of $\qSchu$ can be identified with $\cup_{n\in \N}\Lambda_{\text{st}}^{\ell}(n)$ by deleting the leftmost $\emptyset$.
Let $\Par^\ell(n)$  be the set of $\ell$-multipartitions  of $m$ with usual dominance order $\lhd$. For $\la\in \Par^\ell(n)$, we denote $|\la|=n$. 

Recall for any $\lambda\in\Par^\ell(n)$ a $\lambda$-tableau $\mathbf S$ is obtained from the Young diagram $Y(\lambda)$ by inserting ordered pairs $(i,p)$, also denoted by $i_p$, into $Y(\lambda)$, where $1\le p\le \ell$ and $i$ is a positive integer.
For $\lambda\in\Par^\ell(n)$ and $\mu\in \Lambda_{st}^\ell(n)$, let $\SST(\lambda,\mu)$ be the set of all semistandard $\lambda$-tableaux of type $\mu$ as defined in \cite{DJM98}. Associated to a semistandard tableau $\bT\in \SST(\lambda,\mu)$  we can construct a matrix $A_{\bT}\in \Mat_{\overline{\mu},\overline{\lambda}}$ as in \cite[\S 4.5]{SW24Schur} such that $a_{(p,i),(q,j)}$ is the number of $i_p$ in $j$th row of $\mathbf T^{(q)}$, where $\bar \lambda $ is the strict composition obtained from the multi-composition  $\lambda$ by a forgetful map. Thus, we have a reduced chicken foot ribbon diagram $[A_{\bT}]$  of shape $A_{\bT}$. Let $[\bT]$ denote the ornamentation associated with $(A_{\bT},(\emptyset))$. We refer to \cite[Example 4.8]{SW24Schur} for an explicit example for the $A_\bT$ and the diagram $[\bT]$, where each crossing therein is now chosen to be positive.


The notion of cellular bases was introduced by Graham-Lehrer \cite{GL96}. Let $\bfS_\bfu(n)$ be the path algebra of the subcategory  $\qSchu(n)$ with object set $\Lambda_{\text{st}}^{\ell}(n)$. 
We have the following diagrammatic description of a cellular basis on $\bfS_\bfu(n)$, which is compatible with its counterpart on the cyclotomic $q$-Schur algebra \cite{DJM98} via the algebra isomorphism established in \cite{SSW25}. 

\begin{proposition}  \cite[Theorems 5.12, 5.13]{SSW25}
    The algebra $\bfS_\bfu(n)$ has a cellular  basis given by 
\begin{align}  \label{doubleSST}
\bigcup_{\overset{\lambda\in \Par^\ell(n)}{\mu,\nu\in \Lambda_{\text{st}}^\ell(n)} }
\big\{ [\bT]\circ [\bS]^\div \mid \bT\in \SST(\lambda,\nu), \bS\in \SST(\lambda,\mu) \big\}, 
\end{align}
where the required anti-involution $\div$ is defined by rotating a string diagram by 180 degrees around a horizontal axis.
\end{proposition}

By \cite[Theorem~ 5.17]{SSW25}, we have a fullly faithful functor from $\qW_\bfu$ to $\qSchu$ which sends any diagram $f$ to 
 $\begin{tikzpicture}[baseline = 10pt, scale=0.5, color=\clr]
\draw[-,line width =1pt,color=\cred] (-4.2,.2) to (-4.2,1.6);
\draw (-4.2,0) node{$\scriptstyle \red{u_1}$};    
\draw[-,line width =1pt,color=\cred] (-3.5,.2) to (-3.5,1.6);
\draw (-3.4,0) node{$\scriptstyle \red{u_2}$};
\draw(-2.6,.8) node{$\ldots$};
\draw[-,line width =1pt,color=\cred] (-2,.2) to (-2,1.6);
\draw (-1.9,0) node{$\scriptstyle \red{u_{\ell}}$};
\draw (-1.2,1) node{$f$};
\end{tikzpicture}.
$
The set $\Lambda_\st(n)$ of objects is sent to 
\[
\bar\Lambda^\ell(n):=\{ (\emptyset^{\ell-1}, \mu)\mid \mu\in \Lambda_\st(n)\}.
\]
The above functor induces an isomorphism of algebras 
\[
\phi_n: \wS_\bfu(n) \cong  {\bf1}_n \bfS_\bfu(n){\bf1}_n = \bigoplus_{\nu,\mu\in \bar \Lambda^\ell(n)} \Hom_{\qSchu}(\nu,\mu)  \] 
where  ${\bf1}_n=\oplus_{\nu\in \bar\Lambda^{\ell}(n)}1_\nu$ .
Below we shall identify $\wS_\bfu(n)$ with ${\bf1}_n \bfS_\bfu(n){\bf1}_n$ via the isomorphism $\phi_n$ and construct a cellular basis of $\wS_\bfu(n)$. 

Note that for any $\bT\in \SST(\lambda,\nu), \bS\in \SST(\lambda,\mu)$, 
\begin{equation}
  {\bf1}_n [\bT]\circ [\bS]^\div {\bf1}_n= \begin{cases}
    [\bT]\circ [\bS]^\div, & \text{ if }\mu,\nu \in \bar\Lambda^\ell(n)\\
    0,& \text{otherwise}.
\end{cases} 
\end{equation}
This implies the cellular basis for $\wS_\bfu$ as follows. 
\begin{theorem}
  \label{thm:cellularforwebu}
The algebra $\wS_\bfu(n)$ admits a cellular basis  
\begin{align}  \label{doubleSSTweb}
\bigcup_{\overset{\lambda\in \Par^\ell(n)}{\mu,\nu\in \bar\Lambda^\ell(n)} }
\big\{ [\bT]\circ [\bS]^\div \mid \bT\in \SST(\lambda,\mu), \bS\in \SST(\lambda,\nu) \big\}. 
\end{align}
\end{theorem}
 For any $\lambda\in \Par^\ell(n)$, the (left) cell module   $\Delta^\lambda$ for $\wS_\bfu(n)$ is the module with basis 
 \[ 
 \{[\bT]\circ [\bS]^\div + \wS_\bfu(n)^{\rhd\lambda} \mid \mu\in \bar\Lambda^\ell(n), \bT \in \SST(\lambda,\mu)\} 
 \]
for any fixed $\bS\in \SST(\lambda, \nu)$, where $\wS_\bfu(n)^{\rhd\lambda}$ is the span of elements $[\bT']\circ [\bS']^\div $ in \eqref{doubleSSTweb}
such that $\bT'\in \SST(\lambda',\nu') $, $\bS\in \SST(\lambda',\mu')$ and $\lambda'\rhd \lambda$.
Note that the cell module does not depend on the choice of $\bf S$. So, we even can write
$[\bT]\circ [\bS]^\div + \wS_\bfu(n)^{\rhd\lambda}$ as $[\bT]$ if there is no confusion.
\begin{rem}
\label{rem:cellularforcychecke}
Multiplying both sides of \eqref{doubleSSTweb} by the idempotent $1_{(1^n)}$, we obtain the well-known  cellular basis of the cyclotomic Hecke algebra $\bfH_\bfu(n)$ (e.g., \cite{DJM98}) 
\begin{align}  \label{doubleSSTheck}
\bigcup_{ \lambda\in \Par^\ell(n) }
\big\{ [\bT]\circ [\bS]^\div \mid \bS, \bT\in \SST(\lambda,(\emptyset^{\ell-1},(1^n))  \big\}. 
\end{align}
Then for any $\lambda\in  \Par^\ell(n)$ we have the cell (Specht) module $ S^\lambda$ of $\bfH_{\bfu}(n)$ with respect to the cellular basis above. Obviously, 
\begin{equation}
\label{equ:idemoncell}
 1_{(1^n)} \Delta^\lambda \cong  S^\lambda, \qquad
 \text{ for } \lambda \in  \Par^\ell(n).   
\end{equation}
\end{rem}

\subsection{Cell filtrations on $\Res \Delta^\lambda$ and $\Ind \Delta^\lambda$}

We identify a multi-partition $\lambda=(\lambda^{(1)}, \ldots, \lambda^{(\ell)})\in \Par^\ell(n)$ with its Young diagram $[\lambda]$. We write a node $x\in [\lambda]$ as $ (i,j,k)$ if $x$ is at $i$th row and $j$th column of the $k$-th component. We say $x$ is  removable (resp. addable) of $\lambda$ if $\lambda\backslash x$ (resp. $\lambda\cup x$) is a multi-partition of $n-1$ (resp. $n+1$). 

A set of nodes $\{x_1,x_2,\ldots,x_a\}$ is called removable (resp., addable) if $\lambda\backslash \{x_1,x_2,\ldots,x_a\}$ (resp. $\lambda\cup \{x_1,x_2,\ldots,x_a\}$) is a multi-partition of $n-a$ (resp. $n+a$). Let
\begin{align}
\begin{split}
    \mathcal R(\lambda)_a &=\{ \text{$\ell$-multipartitions  obtained from $\lambda$ by removing $a$ removable nodes}\},
    \\
    \bar {\mathcal R}(\lambda)_a &=\{\mu \in \mathcal R(\lambda)_a\mid \mu =\lambda\backslash \{x_1, \ldots,x_a\} \text{ with all
$x_i$ in distinct columns}\}.
\end{split}
\end{align}

Similarly, we have the set $\mathcal A(\lambda)_a$ of $\ell$-multipartitions obtained by adding $a$ addable nodes to $\lambda$, and 
\begin{align}
    \bar {\mathcal A}(\lambda)_a &=\{\mu \in \mathcal A(\lambda)_a\mid \mu =\lambda\cup \{x_1, \ldots,x_a\} \text{ with all
$x_i$ in distinct columns}\}.
\end{align}
In particular, we have $\bar {\mathcal R}(\lambda)_1=  \mathcal R(\lambda)_1$
and $\bar {\mathcal A}(\lambda)_1=  \mathcal A(\lambda)_1$.

Recall that a   $\lambda$-tableau $\mathbf S$ in $\SST(\lambda, \mu)$ is obtained from Young diagram $Y(\lambda)$ by inserting ordered pairs $(i,p)$, also denoted by $i_p$, into $Y(\lambda)$ satisfying certain conditions,  where $1\le p\le \ell$ and $i$ is a positive integer (cf. \cite[Definition (4.4)]{DJM98} or \cite[Ex. 4.2]{SW24Schur}). In particular, the number of times of $(i,p)$ as an entry in $\mathbf S$ is $\mu^{(p)}_i$, for all $i,p$. Let $\bar \Lambda^\ell(n)_a$ be the subset of $\bar\Lambda^\ell(n)$ consisting of all  elements such that  $\mu^{(\ell)}=(\mu_1,\ldots,\mu_k)$ with  $\mu_k=a$.
For any $\bT\in \SST(\lambda, \mu)$ with $\mu\in \bar\Lambda^{\ell}(n)$ 
we have $\textbf{1}_{a,n}[\bT]\neq 0$ only if $\mu\in \bar \Lambda^\ell(n)_a$.
In this case, we define $\bT^a$ to be the tableaux obtained from $\bT$ by deleting all $a$ of $k_{\ell}$'s. In other words, the diagram $[\bT^a]$ is obtained from $[\bT]$ by deleting the strands that connect to $\mu_k$.  

Since there are $a$ nodes $R=\{x_1,\ldots, x_a\}$ filled by $k_{\ell}$, we see that  $\bT^a\in \SST(\lambda\backslash R, \mu\backslash a ) $, where $\mu\backslash a\in \bar\Lambda^{\ell}(n-a)$ such that $(\mu\backslash a)^{(\ell)}=(\mu_1,\ldots,\mu_{k-1})$. In particular, 
since $\bT$ is semi-standard, $x_i$ and $x_j$ are not in the same column and hence $\lambda\backslash R\in \bar {\mathcal R}(\lambda)_a$. Moreover, any $\bS\in \SST(\lambda\backslash R, \mu\backslash a )$
can be obtained in this way such that $\bS=\bT^a$.

Recall there is a unique $\bT^{\lambda\backslash R}\in \SST(\lambda\backslash R, \lambda\backslash R) $ such that $ [\bT^{\lambda\backslash R}]=1_{\lambda\backslash R}$.
Let $\bT^{\lambda\backslash R}_a \in \SST(\lambda,\lambda\backslash R)$ be the tableaux obtained from $\bT^{\lambda\backslash R}$  by adding the nodes  $R$ back. 
Then we see that 
\[
[\bT]=[\bT^a] \circ [\bT_a^{\lambda\backslash R}]. 
\]
Recall the restriction functor $\Res^n_{n-a}: \wS_\bfu(n)\modf \rightarrow\wS_\bfu(n-a)\modf$ from \eqref{eq:IndRes}.

\begin{proposition}  \label{prop:restriction}
    For any $\lambda\in \Par^\ell(n)$, there exists a filtration of $\wS_\bfu(n-a)$-modules 
    \[0=M_{k+1}\subset M_k\subset \ldots \subset M_1=\Res^n_{n-a} (\Delta^\lambda) \]
    such that $M_i/M_{i+1}\cong \Delta^{\mu^i}$, 
    where $\{\mu^1\lhd \mu^2\lhd\ldots\lhd\mu^k \}=\bar {\mathcal R}(\lambda)_a$. 
\end{proposition}

\begin{proof}
    Define 
    \[M_j=\{[\bT] \mid \bT \in \SST(\lambda,\mu), \mu\in \bar\Lambda^\ell(n)_a,  
    \bT^a\in \SST(\mu^i,\mu\backslash a), i\ge j  \} .\]
    It suffices to show that the linear map 
    \begin{align*}
        \psi : \Delta^{\mu^j} \longrightarrow M_j/M_{j+1}, \qquad
      [\bT^a] \mapsto  [\bT] + M_{j+1},
    \end{align*}
    is a homomorphism (and hence an isomorphism) of $\wS_\bfu(n-a)$-modules.

For any $h\in \wS_\bfu(n-a)$, we have 
\[ h [\bT^a]=\sum_{\bS^a\in \SST(\mu^j,\mu\backslash a)} c_{\bT^a,\bS^a}
[\bS^a] +\sum_{\nu\rhd \mu^j, \bT_1\in \SST(\nu,\gamma),\bS_1\in \SST(\nu,\mu^j)} (*) \,[\bT_1]\circ [\bS_1]^\div \]
It suffices to establish the claim that 
\[([\bT_1]\circ [\bS_1]^\div)\otimes 1_a \circ [\bT_a^{\mu^j}]\in M^{j+1}.\]
 
  Note that $( [\bS_1]^\div\otimes 1_a) \circ [\bT_a^{\mu^j}]\in \Hom_{\qW_\bfu}(\nu\cup a, \lambda)$.
 By \cite[Proposition 5.11]{SSW25}, 
  \[ ( [\bS_1]^\div\otimes 1_a) \circ [\bT_a^{\mu^j}]= \sum_{\bT_2,\bS_2} (*)\,[\bT_2]\circ [\bS_2]^\div   \]
  where the sum is over $\bT_2\in \SST(\tilde \nu,\nu\cup a)$ and $\bS_2\in \SST(\tilde \nu,\lambda)$.
  Note that $ \SST(\tilde \nu,\lambda)\neq \emptyset$ only if $ \lambda \unlhd \tilde\nu$. 
  If   $\tilde \nu \ntrianglelefteq \lambda$, then 
  $[\bT_2]\circ [\bS_2]^\div=0$ in $\Delta^\lambda$ and the claim trivially holds. 
  So, it remains to consider the case 
  \[ \mu^j\cup a \unlhd \nu\cup a \unlhd\tilde \nu = \lambda.\]
  This condition implies $\nu\in \{ \mu^{j+1},\ldots,\mu^k\}$.
  Furthermore, $ ([\bS_1]^\div \otimes 1_a) \circ \bT^{\mu^j}_a =\bT^{\mu^h}_a$ for some $ \nu=\mu^h$ with $j+1\le h\le k$.
  Hence the claim holds. This completes the proof.
\end{proof}

\begin{proposition}
\label{prop:induction}
    For any $\lambda\in \Par^\ell(n)$, there exists a filtration of $\wS_\bfu(n+a)$-modules 
    \[0=M_{k+1}\subset M_k\subset \ldots \subset M_1=\Ind^{n+a}_{n} (\Delta^\lambda) \]
    such that $M_i/M_{i+1}\cong \Delta^{\mu^i}$, 
    where $\{\mu^1\lhd \mu^2\lhd\ldots\lhd\mu^k \}=\bar {\mathcal A}(\lambda)_a$. 
\end{proposition}

\begin{proof}
    By definition, we have 
    \begin{align*}
        \Ind^{n+a}_{n} (\Delta^\lambda)&= \wS_{\bfu}(n+a) \textbf{1}_{a,n+a}\otimes _{\wS_\bfu(n)} \Delta^\lambda\\
        &=\wS_{\bfu}(n+a) \textbf{1}_{a,n+a}\otimes _{\wS_\bfu(n)}  \wS_\bfu(n) 1_\lambda\circ [\bS]^\div / \wS_\bfu^{\rhd\lambda }\\
        &= M^\lambda/N^\lambda,
    \end{align*} 
    where, for any fixed $\bS\in \SST(\lambda, (\emptyset^{\ell-1},\bar \lambda))$ for some $\bar\lambda\in\Lambda_\st(n)$,
   \begin{align*}
       M^\lambda&= \wS_{\bfu}(n+a) \textbf{1}_{a,n+a}\otimes _{\wS_\bfu(n)}  \wS_\bfu(n) 1_\lambda\circ [\bS]^\div, \\  N^\lambda&=\wS_{\bfu}(n+a) \textbf{1}_{a,n+a}\otimes _{\wS_\bfu(n)}\wS_\bfu^{\rhd\lambda }.
   \end{align*} 
    We write $\lambda\cup a=(\lambda^{(1)}, \lambda^{(2)},\ldots ,\lambda^{(\ell)}\cup a)$, where $\lambda^{(\ell)}\cup a $ is obtained by attaching the part $a$ to the end of $\lambda^{(\ell)}$. Using the cellular basis of $\wS_\bfu(n+a)$ in Theorem \ref{thm:cellularforwebu}, we see that 
    $M^\lambda$ is spanned by 
\[\{ [\bT_1]\circ [\bT_2]^\div \otimes [\bS]^\div \mid \bT_1\in \SST(\mu, \nu), \bT_2\in \SST(\mu, \lambda\cup a), \mu\in \Par^\ell(n+a), \nu \in \bar\Lambda^{\ell}(n+a)\}.\]
Let $R=\{x_1,\ldots,x_a\}$ be the nodes in $\bT_2$ filled by $k_\ell$ where $k=l(\lambda^{(\ell)})+1$. 
Recall that  $\bT_2^a=\bT_2\backslash R\in \SST(\mu\backslash R, \lambda)$.  If this set is not $\emptyset$, we must have $ \mu\backslash R\unrhd \lambda$. %or $\mu\backslash R=\lambda$. 

If $ \mu \backslash R\rhd \lambda$, then $ [\bT_2 ]^\div= [\bT_a^{\mu\backslash R}]\div [\bT_2^a]^\div $ 
and hence $ [\bT_1]\circ [\bT_2]^\div \otimes [\bS]^\div=[\bT_1]\circ[\bT_a^{\mu\backslash R}]\div\otimes [\bT_2^a]^\div[\bS]^\div \in N^\lambda  $ since $ \mu \backslash R\rhd \lambda$. 

If $\mu \backslash R =\lambda$, then we must have $\mu\in \bar{\mathcal A}(\lambda)_a$  and $[\bT_2^a]=1_\lambda$ and $\bT_2=\bT_a^{\mu\backslash R} $.
This implies that $ \Ind^{n+a}_{n} (\Delta^\lambda)$ is spanned by 
\[
B^\lambda :=\Big\{[\bT_1]\circ (\bT_a^{\lambda})^\div \otimes [\bS]^\div  +N^\lambda ~\Big |~ \bT_1\in \SST(\mu,\nu),\mu\in  \bar{\mathcal A}(\lambda)_a , \nu \in \bar\Lambda^{\ell}(n+a) \Big\}.
\] 

Note that the dimension of $\Ind^{n+a}_{n} (\Delta^\lambda) $ is independent of the field and parameters by the proof of  Lemma \ref{Lem:Projective}, and so we can assume that $\wS_\bfu$ is seimisimple when we compute its dimension. In this case, using the result on the restriction functor in Proposition~ \ref{prop:restriction} and the fact that $\Ind_n^{n+a}$ is left adjoint to  
$ \Res_n^{n+a}$, we see that this dimension is $|B^\lambda|$ by the Frobenius reciprocity. Therefore $B^\lambda $ is a basis of $\Ind^{n+a}_{n} (\Delta^\lambda)$.

Now we define 
\[
M_j= \Big\{[\bT_1]\circ (\bT_a^{\lambda})^\div \otimes [\bS]^\div  +N^\lambda ~\Big |~\bT_1\in \SST(\mu^i,\nu), i\ge j , \nu \in \bar\Lambda^{\ell}(n+a) \Big\}.
\]
The above arguments also show that  the linear map 
    \begin{align*}
        \psi : \Delta^{\mu^j}&\longrightarrow M_j/M_{j+1}, \qquad
      [\bT^a] \mapsto  [\bT] + M_{j+1}
    \end{align*}
    is a homomorphism (and hence an isomorphism) of $\wS_\bfu(n+a)$-modules.
\end{proof}

\subsection{Functors $\iInd$ and $\iRes$}

We specialize at $v=\eps$, where $\eps$ is either generic (i.e., not a root of 1) or $\eps=\exp(2\pi i/p')$, a primitive $p'$-th root of 1. We set $p=p'$ if $p'$ is odd and $p=p'/2$ if $p'$ is even. In other words, 
\begin{align}
    \text{$\text{order}(\eps^2)=p$   in all cases (including $p=\infty$ for $\eps$ generic). }
\end{align}
Set 
\begin{align} \label{eq:I}
\I =\Z \text{ or } \Z/p\Z, 
    \end{align}
corresponding to the case for $\eps$ generic or when $\text{order}(\eps^2)=p$. We denote $\Z/p\Z =\{\overline{0}, \overline{1}, \ldots, \overline{p-1}\}$. 

Let $\bfs =(s_1, \ldots, s_\ell) \in \I^\ell$ and $\bfu =(u_1,\ldots, u_\ell) \in (\kk^\times)^\ell$. We shall make the assumption on parameters that 
\begin{align} \label{u:parameter}
(u_1,\ldots, u_\ell)= (\eps^{2s_1}, \ldots, \eps^{2s_\ell}), \qquad \text{ or } \bfu =\eps^{2\bfs} \text{ for short}. 
\end{align}
Note this is well defined by \eqref{eq:I}. The arbitrary parameter cases can be reduced to this by Morita equivalences (cf. \cite{DM02}). 
For any node $x=(i,j,k)\in [\lambda]$, we define the residue of $x$ by 
\[
\res(x)=u_k\eps^{-2(j-i)} = \eps^{s_k-2(j-i)}
.\]

It is known that the central elements 
$f(X_1,\ldots,X_n)\in \kk[X_1^\pm,\ldots,X_n^\pm]^{\mathfrak S_n} $ acts on $1_{(1^n)}\Delta^\lambda$ (the corresponding cell module of the cyclotomic Hecke algebra) as the scalar
\[f(\res(x_1), \ldots, \res(x_n))\]
where $x_1,\ldots, x_n$ are all the nodes of $[\lambda]$. 
Let $\gamma_\lambda=(\res(x_1), \ldots, \res(x_n))$. Then 
$\Delta^\lambda\in \wS_\bfu(n)\modf_{\gamma_\lambda}$ and we can identify $\gamma_\lambda$ with $\sum_{j=1}^n\alpha_{t_j}$ if $ \gamma_\lambda=(\eps^{-2t_1}, \ldots, \eps^{-2t_n} )$. Here $\{\alpha_i|i\in \I\}$ denote the simple roots of $\widehat {\mathfrak {sl}}_p$.

Similar to $\gamma$, we can identify the multiset $\bfi=\{i_1,\ldots, i_a\}$ with $ \alpha_{i_1}+\ldots+ \alpha_{i_a}$, and denote the multiset $\eps^{-2\bfi}:=\{\eps^{-2i_1},\ldots, \eps^{-2i_a}\}$.  We denote $|\bfi|=a$ and let 
\begin{align*}
    \bar{ \mathcal R}(\lambda)_{\bfi} &=\Big\{\mu=\lambda\backslash\{x_1,\ldots,x_a\} \in \bar {\mathcal R}(\lambda)_{|\bfi|} ~\Big |~ 
    \{\res(x_1),\ldots, \res(x_a) \}= \eps^{-2\bfi} 
    \Big\}.
\end{align*}
Similarly, we have the subset 
$\bar{ \mathcal A}(\lambda)_{\bfi} $ of $\bar{ \mathcal A}(\lambda)_{|\bfi|}$.

Similar to $\bfc\text{-}\Res^n_{n-|\bfk|}$ from \eqref{iRes_affine_n} in the affine setting (now with $\bfc=\eps^{-2\bfi}$), we define a functor $\eps^{-2\bfi}\text{-}\Res^n_{n-|\bfi|}$ (and simply denoted by $\iRes^n_{n-|\bfi|}$ from now on) sending any $M\in \wS_\bfu(n)\modf_\gamma$ to $\text{Proj}_{\gamma\setminus\eps^{-2\bfi}}\circ \text{Res}^n_{n-|\bfi|}(M) \in \wS_\bfu(n-|\bfi|)\modf_{\gamma\setminus\eps^{-2\bfi}} $.
%, where $\text{Proj}_{\gamma\setminus\eps^{-2\bfi}}$ is the projection from $ \wS_\bfu(n-|\bfi|)\modf$ to $\wS_\bfu(n-|\bfi|)\modf_{\gamma\setminus\eps^{-2\bfi}}$. 
Similarly, we have a functor $\iInd^{n+|\bfi|}_n$ (which is a simplified notation for $\eps^{-2\bfi}\text{-}\Ind^{n+|\bfi|}_n$), which sends $M\in \wS_\bfu(n)\modf_\gamma$ to a module in $\wS_\bfu(n+a)\modf_{\gamma\cup\eps^{-2\bfi}} $. Similarly, we will denote the linear operators $f_\bfc$ with $\bfc={\eps^{-2\bfi}}$ by $f_\bfi$. The following result follows from  Propositions \ref{prop:restriction}--\ref{prop:induction}.

\begin{proposition}
\label{prop:IndRes}
Let $\bfi \in \Z \I$. 
\begin{enumerate}
    \item     For any $\lambda\in \Par^\ell(n)$, there exists a filtration of $\wS_\bfu(n-|\bfi|)$-modules 
    \[0=M_{k+1}\subset M_k\subset \ldots \subset M_1=\iRes^n_{n-|\bfi|} (\Delta^\lambda) \]
    such that $M_i/M_{i+1}\cong \Delta^{\mu_i}$, 
    where $\{\mu_1\lhd \mu_2\lhd\ldots\lhd\mu_k \}=\bar {\mathcal R}(\lambda)_\bfi$.
    \item For any $\lambda\in \Par^\ell(n)$, there exists a filtration of $\wS_\bfu(n+|\bfi|)$-modules 
    \[0=M_{k+1}\subset M_k\subset \ldots \subset M_1=\iInd^{n+|\bfi|}_{n} (\Delta^\lambda) \]
    such that $M_i/M_{i+1}\cong \Delta^{\mu_i}$, 
    where $\{\mu_1\lhd \mu_2\lhd\ldots\lhd\mu_k \}=\bar {\mathcal A}(\lambda)_\bfi$. 
\end{enumerate}    
\end{proposition}

We then define functors
\[
\iRes:= \oplus_{n\in \N}\iRes^n_{n-|\bfi|}, \qquad \iInd:= \oplus_{n\in\N}\iInd^{n+|\bfi|}_n
\] 
on $\qW_\bfu\modf.$


\subsection{A semisimplity criterion}


\begin{proposition}
\label{prop:semisimple} 
Suppose $\bfu=(u_1,\ldots,u_\ell)\in (\kk^*)^\ell$. Then $\wS_\bfu(n)$ is semisimple if and only if the cyclotomic Hecke algebra $\bfH_{\bfu}(n)$ is semisimple. 
\end{proposition}

\begin{proof}
   The only if part holds  since $\bfH_{\bfu}(n)= 1_{(1^n)}\wS_{\bfu}(n) 1_{(1^n)}$.

   If $\bfH_{\bfu}(n)$ is semisimple, then $ \bfS_\bfu(n)$ is semisimple since the latter is the endomorphism algebra of a direct sum of permutation modules of $\bfH_{\bfu}(n)$. Now $\wS_{\bfu}(n)$ is semisimple since 
   $ \wS_{\bfu}(n)=e\bfS_\bfu(n)e$, where the idempotent $e=\sum_\mu 1_\mu$ sums over all $\ell$-multicompositions $\mu$ of $n$ such that $\mu^{(k)}=\emptyset$ for $k<\ell$.
\end{proof}

The criterion for semisimplicity of the cyclotomic Hecke algebra $\bfH_{\bfu}(n)$ was worked out by Ariki; see \cite{Ar02}.
